\documentclass[11pt]{article}
\usepackage[margin=18mm,headheight=22pt]{geometry}
\usepackage{amsmath}
\usepackage{fontspec}
\usepackage{unicode-math}
\setmainfont{NewCM10}[Extension=.otf,UprightFont=*-Book,BoldFont=*-Bold,ItalicFont=*-BookItalic,BoldItalicFont=*-BoldItalic]
\setsansfont{NewCMSans10}[Extension=.otf,UprightFont=*-Book,BoldFont=*-Bold,ItalicFont=*-BookOblique,BoldItalicFont=*-BoldOblique]
\setmathfont{NewCMMath-Book.otf}
\usepackage{xcolor}
\definecolor{FAUblue}{HTML}{04316A}
\definecolor{Paperblue}{HTML}{EEF2F7}
\definecolor{EEIgold}{HTML}{CC9C42}
\usepackage{etoolbox,enumitem,booktabs,tabularx,array,needspace,microtype}
\usepackage[most]{tcolorbox}
\usepackage{titlesec,fancyhdr}
\usepackage[hidelinks]{hyperref}
\usepackage{draftwatermark}
\SetWatermarkText{Made by Thanate Amornjiraporn (Focus)}
\SetWatermarkColor{black!7}\SetWatermarkFontSize{30pt}\SetWatermarkAngle{45}
\setlength{\parindent}{0pt}
\setlength{\parskip}{5pt}
\setlist{itemsep=2pt,topsep=3pt,leftmargin=1.6em,beginpenalty=10000,endpenalty=10000}
\raggedbottom
\clubpenalty=10000 \widowpenalty=10000 \displaywidowpenalty=10000 \interlinepenalty=10000
\setcounter{secnumdepth}{0}
\setcounter{tocdepth}{1}
\pagestyle{fancy}\fancyhf{}
\fancyhead[L]{\small\sffamily\color{FAUblue} From Fundamentals to AI Apocalypse · APDR · Exercise 0}
\fancyhead[R]{\small\sffamily Computer Exercise 0: Practice solutions}
\fancyfoot[C]{\small\thepage}
\titleformat{\section}{\Large\bfseries\sffamily\color{FAUblue}}{}{0pt}{}
\titleformat{\subsection}{\large\bfseries\sffamily\color{FAUblue}}{}{0pt}{}
\titlespacing*{\section}{0pt}{16pt}{6pt}
\preto\section{\Needspace{8\baselineskip}}
\preto\subsection{\Needspace{6\baselineskip}}
\newtcolorbox{equationcard}[1]{colback=Paperblue,colframe=Paperblue,boxrule=0pt,arc=0pt,left=10pt,right=10pt,top=6pt,bottom=6pt,title={#1},fonttitle=\small\bfseries\sffamily,coltitle=FAUblue,before skip=8pt,after skip=8pt}
\newtcolorbox{workedbox}[1]{colback=white,colframe=EEIgold,boxrule=0pt,leftrule=3pt,arc=0pt,left=10pt,top=4pt,bottom=4pt,title={#1},fonttitle=\small\bfseries\sffamily,coltitle=FAUblue,attach title to upper={\par},before skip=8pt,after skip=8pt}
\newtcolorbox{recallbox}[1]{colback=white,colframe=FAUblue,boxrule=0pt,leftrule=3pt,arc=0pt,left=10pt,top=4pt,bottom=4pt,title={#1},fonttitle=\small\bfseries\sffamily,coltitle=FAUblue,attach title to upper={\par},before skip=8pt,after skip=8pt}
\usepackage{environ}
\newsavebox\measurebox
% Typeset the section once invisibly to measure it; start a new page if it will not fit in what is left.
\makeatletter
\newif\ifkeep@measuring
\NewEnviron{keepsection}{%
  \ifkeep@measuring \BODY \else
    \edef\keep@sec{\the\value{section}}\edef\keep@sub{\the\value{subsection}}%
    \setbox\measurebox\vbox{\hsize\linewidth\keep@measuringtrue\def\Needspace##1{}\BODY}%
    \setcounter{section}{\keep@sec}\setcounter{subsection}{\keep@sub}%
    \dimen2=\ht\measurebox \advance\dimen2\dp\measurebox
    \dimen@=\textheight
    \ifdim\pagegoal<\maxdimen
      \dimen@=\pagegoal \advance\dimen@-\pagetotal
      \ifdim\dimen2>\dimen@ \clearpage \dimen@=\textheight \fi
    \fi
    % if it fits, nothing inside may force an early break
    \ifdim\dimen2>\dimen@ \BODY \else \begingroup\def\Needspace##1{}\BODY\endgroup \fi
  \fi}
\makeatother
\newtcolorbox{notebox}{colback=Paperblue,colframe=Paperblue,boxrule=0pt,arc=0pt,left=8pt,right=8pt,top=4pt,bottom=4pt,before skip=6pt,after skip=6pt}
\newtcolorbox{plainbox}{colback=white,colframe=FAUblue,boxrule=0pt,leftrule=3pt,arc=0pt,left=10pt,right=8pt,top=4pt,bottom=4pt,before skip=6pt,after skip=8pt}
\usetikzlibrary{arrows.meta}
\title{\sffamily\bfseries\color{FAUblue}Computer Exercise 0: Practice solutions}
\author{Algorithms, Programming, and Data Representation \ensuremath{\cdot} Module 92313 \ensuremath{\cdot} Exercise 0\\[2pt]\small From Fundamentals to AI Apocalypse · personal study guide}
\date{}
\begin{document}
\maketitle
\begingroup\footnotesize\setlength{\parskip}{0pt}\tableofcontents\endgroup
\clearpage

\Needspace{5\baselineskip}
Practice \ensuremath{\cdot} not course-graded

\Needspace{8\baselineskip}
{\large Exercise 0 is a practice workflow check. Use it to verify your submission setup; it is not part of the course grade. Try each task before opening the worked steps.}


\begin{keepsection}
\Needspace{5\baselineskip}
\section{Task 1: Return the exact greeting}

\textbf{Connected topic:} Computers, programs, and data representation (Chapter 2)

\begin{keepsection}
\Needspace{6\baselineskip}
\subsection{Task 01 \ensuremath{\cdot} hello()}

\Needspace{3\baselineskip}
Try first: write the function first, then compare its returned string character for character.

\nopagebreak[4]{\small\itshape Answers: see “Task 1: Return the exact greeting” under Answers at the end.}
\end{keepsection}
\end{keepsection}

\begin{keepsection}
\Needspace{5\baselineskip}
\section{Task 2: Add two numbers}

\textbf{Connected topic:} Computers, programs, and data representation (Chapter 2)

\begin{keepsection}
\Needspace{7\baselineskip}
\subsection{Task 02 \ensuremath{\cdot} add\_two\_numbers(a, b)}

\Needspace{4\baselineskip}
Try first: identify the return type implied by the examples and consider negative and decimal inputs.

\nopagebreak[4]{\small\itshape Answers: see “Task 2: Add two numbers” under Answers at the end.}
\end{keepsection}
\end{keepsection}

\begin{keepsection}
\Needspace{5\baselineskip}
\section{Task 3: Square a number}

\textbf{Connected topic:} Computers, programs, and data representation (Chapter 2)

\begin{keepsection}
\Needspace{6\baselineskip}
\subsection{Task 03 \ensuremath{\cdot} square(number)}

\Needspace{3\baselineskip}
Try first: express “square” as multiplication and check a negative input.

\nopagebreak[4]{\small\itshape Answers: see “Task 3: Square a number” under Answers at the end.}
\end{keepsection}
\end{keepsection}

\begin{keepsection}
\Needspace{5\baselineskip}
\section{Task 4: Test whether a number is even}

\textbf{Connected topic:} Computers, programs, and data representation (Chapter 2)

\begin{keepsection}
\Needspace{6\baselineskip}
\subsection{Task 04 \ensuremath{\cdot} is\_even(number)}

\Needspace{3\baselineskip}
Try first: use the remainder after division by two.

\nopagebreak[4]{\small\itshape Answers: see “Task 4: Test whether a number is even” under Answers at the end.}
\end{keepsection}
\end{keepsection}

\begin{keepsection}
\Needspace{5\baselineskip}
\section{Task 5: Build a personalized greeting}

\textbf{Connected topic:} Computers, programs, and data representation (Chapter 2)

\begin{keepsection}
\Needspace{6\baselineskip}
\subsection{Task 05 \ensuremath{\cdot} greet(name)}

\Needspace{3\baselineskip}
Try first: use string formatting and preserve the supplied name exactly.

\nopagebreak[4]{\small\itshape Answers: see “Task 5: Build a personalized greeting” under Answers at the end.}
\end{keepsection}
\end{keepsection}


\clearpage
\section{Answers}
\begin{keepsection}\subsection{Task 1: Return the exact greeting}

\begin{workedbox}{Task 01 \ensuremath{\cdot} hello() \ensuremath{\cdot} worked solution}
\Needspace{10\baselineskip}
\begin{enumerate}
\item The contract is a zero-argument function returning one exact string. There is no input to transform.
\item Return the required literal; do not \texttt{print} it, because printing produces \texttt{None} as the function result.
\item Call the function and compare the value, including capitalization, comma, and exclamation mark.
\item \Needspace{3\baselineskip}
\textbf{The whole function:}
\begin{tcolorbox}[colback=Paperblue,colframe=Paperblue,boxrule=0pt,arc=0pt,left=7pt,right=7pt,top=5pt,bottom=5pt]\small\sloppy
\noindent\hspace*{0em}\texttt{def hello():}\par
\noindent\hspace*{1em}\texttt{return "Hello, World!"}\par
\end{tcolorbox}
\end{enumerate}
\end{workedbox}

\end{keepsection}
\begin{keepsection}\subsection{Task 2: Add two numbers}

\begin{workedbox}{Task 02 \ensuremath{\cdot} add\_two\_numbers(a, b) \ensuremath{\cdot} worked solution}
\Needspace{8\baselineskip}
\begin{enumerate}
\item The required operation is ordinary addition. Keep the parameters generic so Python preserves its numeric behavior for integers and floats.
\item Return \texttt{a + b}; avoid converting to \texttt{int}, which would discard a fractional part.
\item Check that \texttt{add\_two\_numbers(2, 3)} returns 5, then check a negative value and a float.
\item \Needspace{3\baselineskip}
\textbf{The whole function:}
\begin{tcolorbox}[colback=Paperblue,colframe=Paperblue,boxrule=0pt,arc=0pt,left=7pt,right=7pt,top=5pt,bottom=5pt]\small\sloppy
\noindent\hspace*{0em}\texttt{def add\_two\_numbers(a, b):}\par
\noindent\hspace*{1em}\texttt{return a + b}\par
\end{tcolorbox}
\end{enumerate}
\end{workedbox}

\end{keepsection}
\begin{keepsection}\subsection{Task 3: Square a number}

\begin{workedbox}{Task 03 \ensuremath{\cdot} square(number) \ensuremath{\cdot} worked solution}
\Needspace{8\baselineskip}
\begin{enumerate}
\item A square is the number multiplied by itself; it is not the square root and it is nonnegative for real inputs.
\item Return \texttt{number * number}.
\item Check \texttt{square(4)} = 16 and \texttt{square(-3)} = 9.
\item \Needspace{3\baselineskip}
\textbf{The whole function:}
\begin{tcolorbox}[colback=Paperblue,colframe=Paperblue,boxrule=0pt,arc=0pt,left=7pt,right=7pt,top=5pt,bottom=5pt]\small\sloppy
\noindent\hspace*{0em}\texttt{def square(number):}\par
\noindent\hspace*{1em}\texttt{return number * number}\par
\end{tcolorbox}
\end{enumerate}
\end{workedbox}

\end{keepsection}
\begin{keepsection}\subsection{Task 4: Test whether a number is even}

\begin{workedbox}{Task 04 \ensuremath{\cdot} is\_even(number) \ensuremath{\cdot} worked solution}
\Needspace{8\baselineskip}
\begin{enumerate}
\item An integer is even exactly when its remainder modulo 2 is zero. Python’s modulo rule also handles negative integers consistently.
\item Return the Boolean comparison \texttt{number \% 2 == 0}. Do not return the strings \texttt{"\texttt{True}"} or \texttt{"\texttt{False}"}.
\item Check 4 gives \texttt{True}, 5 gives \texttt{False}, and -2 gives \texttt{True}.
\item \Needspace{3\baselineskip}
\textbf{The whole function:}
\begin{tcolorbox}[colback=Paperblue,colframe=Paperblue,boxrule=0pt,arc=0pt,left=7pt,right=7pt,top=5pt,bottom=5pt]\small\sloppy
\noindent\hspace*{0em}\texttt{def is\_even(number):}\par
\noindent\hspace*{1em}\texttt{return number \% 2 == 0}\par
\end{tcolorbox}
\end{enumerate}
\end{workedbox}

\end{keepsection}
\begin{keepsection}\subsection{Task 5: Build a personalized greeting}

\begin{workedbox}{Task 05 \ensuremath{\cdot} greet(name) \ensuremath{\cdot} worked solution}
\Needspace{8\baselineskip}
\begin{enumerate}
\item The output has a fixed prefix, the name, and a final exclamation mark. No punctuation is added to the name itself.
\item Use an f-string to construct and return the full text.
\item Check \texttt{greet("Ada")} returns exactly \texttt{"Hello, Ada!"} and does not \texttt{print} to the terminal.
\item \Needspace{3\baselineskip}
\textbf{The whole function:}
\begin{tcolorbox}[colback=Paperblue,colframe=Paperblue,boxrule=0pt,arc=0pt,left=7pt,right=7pt,top=5pt,bottom=5pt]\small\sloppy
\noindent\hspace*{0em}\texttt{def greet(name):}\par
\noindent\hspace*{1em}\texttt{return f"Hello, \{name\}!"}\par
\end{tcolorbox}
\end{enumerate}
\end{workedbox}

\end{keepsection}

\vfill{\small Personal worked-solution guide based on the APDR Computer Exercise 0 task sheet, rewritten in my own words. It is not an official FAU publication. Exercise 0 is practice; Exercises 1–2 are scored tasks. Use the samples and dataset from your own assigned repository where values are personalized. Made by \textbf{Thanate Amornjiraporn (Focus)}.}

\end{document}
